\documentclass[12pt,a4paper]{article}
\usepackage[T2A]{fontenc}
\usepackage[utf8]{inputenc}
\usepackage[russian]{babel}
\usepackage{amsmath}
\usepackage{geometry}
\usepackage{fancyvrb}
\usepackage{graphicx}
\geometry{left=2cm,right=2cm,top=2cm,bottom=2cm}
\begin{document}
\pagestyle{empty}
\begin{center}
    
    {\fontsize{12pt}{14pt}\selectfont
        МИНИСТЕРСТВО НАУКИ И ВЫСШЕГО ОБРАЗОВАНИЯ\\
        РОССИЙСКОЙ ФЕДЕРАЦИИ\\
        Федеральное государственное бюджетное образовательное учреждение\\
        высшего образования\\
        \textbf{АДЫГЕЙСКИЙ ГОСУДАРСТВЕННЫЙ УНИВЕРСИТЕТ}\\
    }
    \vspace{0.3cm}
    {\fontsize{12pt}{14pt}\selectfont
        НОК "Институт точных наук и цифровых технологий"\\
        Кафедра автоматизированных систем обработки информации и\\
        управления\\
    }
\end{center}
\vspace{4.5cm}
\begin{center}
    
    {\fontsize{14pt}{16pt}\selectfont ОТЧЕТ ПО ПРАКТИКЕ\\}
    \vspace{0.5cm}
    {\fontsize{16pt}{18pt}\selectfont \textbf{Вычисление определителя матрицы произвольного порядка методом Гаусса}\\}
    {\fontsize{14pt}{16pt}\selectfont \textit{Текст из задания по варианту 7}\\}
    \vspace{0.8cm}
    {\fontsize{12pt}{14pt}\selectfont 1 курс, группа УТС\\}
\end{center}
\vspace{3cm}
\hfill
\begin{minipage}{0.48\textwidth}
    
    {\normalsize Выполнил:\\
    \underline{\hspace{3.5cm}} Д. Н. Воробьев\\
    «\underline{\hspace{0.8cm}}» \underline{\hspace{2.2cm}} 2026 г.\\
    \vspace{0.6cm}
    Руководитель:\\
    \underline{\hspace{3.5cm}} С. В. Теплоухов\\
    «\underline{\hspace{0.8cm}}» \underline{\hspace{2.2cm}} 2026 г.}
\end{minipage}
\vfill 
\begin{center}
    {\normalsize Майкоп, 2026 г.}
\end{center}
\newpage

% --- НАЧАЛО ОСНОВНОГО ТЕКСТА ПОСЛЕ ТИТУЛЬНОГО ЛИСТА ---
\pagestyle{plain} % Включаем стандартную нумерацию страниц для последующих листов
\setcounter{page}{2} % Начинаем нумерацию со второй страницы

\begin{center}
    \large\textbf{ВАРИАНТ 7}
\end{center}

\vspace{0.3cm}

\subsection*{Задание}
Найти определитель матрицы.

\vspace{0.3cm}

\subsection*{Теория}

Определителем квадратной матрицы второго порядка $A = \begin{pmatrix} a_{11} & a_{12} \\ a_{21} & a_{22} \end{pmatrix}$ называется число
\[
|A| = a_{11}a_{22} - a_{12}a_{21}.
\]

Определителем $A = \begin{pmatrix} a_{11} & \cdots & a_{1n} \\ \vdots & \ddots & \vdots \\ a_{n1} & \cdots & a_{nn} \end{pmatrix}$ квадратной матрицы порядка $n$, $n \ge 3$, называется число $|A| = \sum_{k=1}^{n} (-1)^{k+1} a_{1k} M_k$, где $M_k$ --- определитель матрицы порядка $n - 1$, полученной из матрицы $A$ вычеркиванием первой строки и столбца с номером $k$.

\vspace{0.5cm}
\hrule
\vspace{0.5cm}

\subsection*{Теоретическая часть}

Определитель (детерминант) матрицы --- это одна из важнейших числовых характеристик квадратной матрицы, которая используется при исследовании и решении систем линейных уравнений, нахождении обратной матрицы, а также в задачах геометрии и механики.

\subsubsection*{1. Понятие минора и алгебраического дополнения}

Для вычисления определителей высших порядков ($n \ge 3$) вводятся понятия миноров и алгебраических дополнений:
\begin{itemize}
    \item \textbf{Минором} $M_{ij}$ элемента $a_{ij}$ квадратной матрицы $A$ порядка $n$ называется определитель матрицы порядка $n-1$, полученной из исходной матрицы путем вычеркивания $i$-й строки и $j$-го столбца.
    \item \textbf{Алгебраическим дополнением} $A_{ij}$ элемента $a_{ij}$ называется его минор, умноженный на коэффициент $(-1)^{i+j}$:
    \[
    A_{ij} = (-1)^{i+j} M_{ij}.
    \]
\end{itemize}

\subsubsection*{2. Теорема о разложении определителя}

Определитель квадратной матрицы равен сумме произведений элементов любой строки (или столбца) на их соответствующие алгебраические дополнения.

Разложение определителя по элементам $i$-й строки записывается как:
\[
|A| = a_{i1} A_{i1} + a_{i2} A_{i2} + \dots + a_{in} A_{in} = \sum_{j=1}^{n} a_{ij} A_{ij}.
\]
Формула общего определения, приведенная в разделе <<Теория>>, представляет собой частный случай этой теоремы --- разложение матрицы по ее первой строке ($i=1$), где слагаемые расписаны через миноры $M_k$ (в обозначениях картинки $M_k = M_{1k}$).

\subsubsection*{3. Ключевые свойства определителей}
\begin{enumerate}
    \item \textbf{Равноправность строк и столбцов:} При транспонировании матрицы её определитель не меняется: $\det(A^T) = \det(A)$.
    \item \textbf{Свойство антисимметрии:} При перестановке двух строк (или столбцов) знак определителя меняется на противоположный.
    \item \textbf{Линейность:} Если все элементы некоторой строки (или столбца) умножить на число $\alpha$, то и весь определитель умножится на это число.
    \item \textbf{Равенство нулю:} Определитель равен нулю, если матрица содержит нулевую строку или две одинаковые (пропорциональные) строки.
    \item \textbf{Инвариантность:} Определитель матрицы не изменится, если к элементам одной строки прибавить элементы другой строки, предварительно умноженные на любое число.
\end{enumerate}
\newpage
\section{Исходный код программы}
Ниже представлен полный исходный код разработанной программы:
\begin{Verbatim}[frame=single, numbers=left, tabsize=4]
#include <iostream>
#include <vector>
#include <iomanip>
#include <limits>
#include <cmath>
#include <clocale>
#include <cstdlib> 
using namespace std;
int main()
{
    system("chcp 65001 > nul");
    int n;
    while (true)
    {
        cout << "Введите размер квадратной матрицы: ";
        if (!(cin >> n))
        {
            cout << "Ошибка! Введите целое число.\n";
            cin.clear();
            cin.ignore(numeric_limits<streamsize>::max(), '\n');
            continue;}
        if (n <= 0)
        {
            cout << "Ошибка! Размер матрицы должен быть больше 0.\n";
            continue;
        }
        break;}
    vector<vector<double>> matrix(n, vector<double>(n));
    cout << "\nВведите элементы матрицы:\n";
    for (int i = 0; i < n; i++)
    {
        for (int j = 0; j < n; j++)
        {
            while (true)
            {
                cout << "A[" << i + 1 << "][" << j + 1 << "] = ";

                if (!(cin >> matrix[i][j]))
                {
                    cout << "Ошибка! Введите число.\n";
                    cin.clear();
                    cin.ignore(numeric_limits<streamsize>::max(), '\n');
                    continue;
                }

                break;
            }
        }
    }
    double determinant = 1.0;
    int swaps = 0;
    for (int i = 0; i < n; i++)
    {
        int pivot = i;
        for (int j = i + 1; j < n; j++)
        {
            if (fabs(matrix[j][i]) > fabs(matrix[pivot][i]))
            {
                pivot = j;
            }
        }
        if (fabs(matrix[pivot][i]) < 1e-10)
        {
            determinant = 0;
            break;
        }
        if (pivot != i)
        {
            swap(matrix[i], matrix[pivot]);
            swaps++;
        }
        for (int j = i + 1; j < n; j++)
        {
            double factor = matrix[j][i] / matrix[i][i];

            for (int k = i; k < n; k++)
            {
                matrix[j][k] -= factor * matrix[i][k];
            }
        }
    }
    if (determinant != 0)
    {
        for (int i = 0; i < n; i++)
        {
            determinant *= matrix[i][i];
        }
        if (swaps % 2 != 0)
        {
            determinant *= -1;
        }
    }
    cout << fixed << setprecision(4);
    cout << "\nОпределитель матрицы = " << determinant << endl;
    return 0;
}

\end{Verbatim}
\newpage
\section{Результат работы программы}
\begin{figure}[h!]
\centering
\includegraphics[widht=0.8\textwidth]{2026-06-22_23-12-15.png}
\caption{Скриншот выполнения программы}
\label{fig:screenshot}
\end{figure}
\end{document}
\end{document}
